\documentclass[11pt,a4paper]{report}

\usepackage[portuguese]{babel}
\usepackage[latin1,utf8]{inputenc}
\usepackage{color}
\usepackage{enumerate}
\usepackage{amsmath}
\usepackage{amssymb}
\usepackage{multirow}
\usepackage{multicol}

\title{\bf Teste de Produção de Relatórios}
\author{Lafaiete Santos}
\date{número de aluno \\ \bf 00000}

\begin{document}

\maketitle

\chapter{Modo Texto}
\vskip 3 mm

\section{Texto}
\noindent Sejam
$$
A = \left\{ x \in R : \int_{0}^{x} f(t)\,dt = 0 \right\}
\quad\text{e}$$
$$
B = \left\{x \in R :
\begin{cases}
F(x) = 0\\
F'(x) = f(x)
\end{cases}
\right\}$$
\noindent onde $f$ é contínua e em $R$ e $F$ é uma primitiva de $f$ tal que $F(0)=0$.\\
\noindent Mostre que $A = B$, indicando explicitamente as inclusões $A \subseteq B$ e $B \subseteq A$.

\section{Listas}
\vskip 2 mm

\begin{flushright}
\Large \bf Revisão de Estrutura de Dados
\end{flushright}

\begin{description}
\item[Conjunto $S$:] Uma coleção de elementos distintos onde a ordem não importa. Exemplo: $\{a,b,c\}$.
\item[Subconjunto ($A \subset B$):] Uma relação onde todos os elementos de $A$ pertencem a $B$.
\end{description}

\begin{enumerate}[I.]
\item \textbf{Considere os conjuntos $A=\{1,2,3\}$ e $B=\{1,2,3,4,5\}$. Analise as afirmações abaixo e selecione a correta:}
\begin{itemize}
\item[$\square$]O elemento $5 \in A$.
\item[$\triangle$]A relação de inclusão correta é $B \subset A$.
\item[$\triangledown$]$A \subset B$, pois todo elemento de $A$ está em $B$.
\item[$\diamondsuit$]$A$ e $B$ são conjuntos disjuntos.
\end{itemize} 
\end{enumerate}

\newpage

\section{Formatações}

\begin{multicols}{2}
\noindent
{\Large \textcolor{blue}{A comunicação científica}} exige clareza e organização.\\
\textcolor{red}{\large Conceitos importantes} devem ser apresentados de forma estruturada, permitindo ao leitor compreender cada ideia.\\
O uso de \textcolor{green}{\small diferentes tamanos de letra} ajuda a destacar partes essenciais do texto e melhora a leitura.

\columnbreak
\noindent {\LARGE \textcolor{magenta}{Destaques visuais}} podem orientar o olhar para elementos relevantes.\\
\textcolor{cyan}{\large A cor}, combinada com espaçamento adequado, torna o documento mais apelativo.\\
Com prática, os alunos aprendem a criar textos \textcolor{blue}{\small claros, organizados e eficazes}.
\end{multicols}

\section{Tabela e referências cruzadas}

\begin{table}[h!]
\centering
\begin{tabular}{|l|c|r|}
\hline
\textbf{Grupo} & \textbf{Expressão} & \textbf{Valor}\\
\hline
\multirow{2}{*}{Cálculo} & $\lim_{x \to 0}\frac{\sin x}{x}$ & $1$
\\ \cline{2-3}
& $\int_{0}^{1} e^{x^2} dx$ & $1.46$ \\ \hline

\multirow{2}{*}{Cat. A} & $\sqrt{\mu^2 + \sigma^2}$ & $12.4$ \\ \cline{2-3}
& $\sum_{i=1}^{n} x_i$ & $10.8$ \\ \hline
\end{tabular}
\label{a}
\end{table}
Como se pode ver na tabela da página \pageref{a}\,...

\section{Outros}

\begin{verbatim}
{\Large \textcolor{blue}{A comunicação científica}} exige clareza e
organização.\\
\textcolor{red}{\large Conceitos importantes} devem ser apresentados
de forma estruturada, permitindo ao leitor compreender cada ideia.\\
O uso de \textcolor{green}{\small diferentes tamanos de letra}
ajuda a destacar parte essenciais do texto e melhora a leitura.
\end{verbatim}

\chapter{Modo Matemático}

\section{Expressão}
$$ 3x + \sqrt{\left(\dfrac{y+3x^2}{5x}\right)^2} + \left[ \dfrac{\frac{2}{3}}{\sqrt[3]{10}} \times \cos(10) \right] $$

\section{Equações}
\begin{gather}
\dfrac{d^2 y}{dx^2} + P(x) \dfrac{dy}{dx} + Q(x)y = 0 \\
f(x) = \sum_{n=0}^{\infty}\dfrac{f^{(n)}(a)}{n!}(x-a)^n
\end{gather}

\section{Sistema de Equações e Matrizes}
$$\Phi(M) =
\begin{cases}

\sum_{i=1}^{n} \lambda_i, & \text{se } M =
\begin{pmatrix}
a_{11} & \cdots & a_{1n} \\
\vdots & \ddots & \vdots \\
a_{n1} & \cdots & a_{nn} \\
\end{pmatrix} \text{e } det(M) \neq 0 \\

\frac{1}{tr(M)}
\begin{bmatrix}
0 & 1\\
1 & 0
\end{bmatrix}, & \text{caso contrário}
\end{cases}
$$

\subsection{Somatório}
$$
S = \sum_{i=1}^{n} \left[\frac{1}{i(i+1)}\sum_{j=1}^{i}(3j^2 + 3j + 1) \right] - \sum_{k=1}^{n} \left(\sum_{m=1}^{k}\dfrac{6m}{k+1}\right)$$

\subsection{Integrais}
\begin{align}
\int_{0}^{1}\left(x^2 + \frac{dF}{dx}\right) dx 
&= \int_{0}^{1} x^2 + F'(x)\,dx \Leftrightarrow \\
\Leftrightarrow \int_{0}^{1} x^2 \, dx + \int_{0}^{1} F'(x)\,dx
&= \frac{1}{3} + F(1) - F(0)
\end{align}

\subsection{Combinações}
$$(x-2)^5= \binom{5}{0}\,x^5+\binom{5}{1}x^{5-1}(-2)^1+\binom{5}{2}x^{5-2}(-2)^2+\cdots+\binom{5}{5}(-2)^5.$$

\end{document}

